\chapter*{Sources, crédits et droits}
\addcontentsline{toc}{chapter}{Sources, crédits et droits}
\begingroup\small
M. Artin, A. Grothendieck, J.-L. Verdier, Théorie des topos et cohomologie étale des schémas (SGA 4), Séminaire de géométrie algébrique du Bois Marie, 1963--1964.\par\medskip
Tome 1 : Théorie des topos. Lecture Notes in Mathematics 269. Springer-Verlag, Berlin--Heidelberg--New York, 1972. Exposés I--IV.\par\medskip
Tome 2 : Théorie des topos. Lecture Notes in Mathematics 270. Springer-Verlag, Berlin--Heidelberg--New York, 1972. Exposés V--VIII, dont V bis.\par\medskip
Tome 3 : Théorie des topos et cohomologie étale des schémas. Lecture Notes in Mathematics 305. Springer-Verlag, Berlin--Heidelberg--New York, 1973. Exposés IX--XIX.\par\medskip
Avec la participation de N. Bourbaki, P. Deligne et B. Saint-Donat.\par\medskip
L'édition LaTeX de Fabrice Orgogozo, état 71766d9, est la base de toutes les couches de ces livres. Ses corrections consignées par la commande sisi et ses notes N.D.E. sont un travail éditorial humain et lui restent attribuées.\par\medskip
Les auteurs et collaborateurs du séminaire sont crédités dans les exposés. Les notices du projet mentionnent notamment Michael Artin, Alexander Grothendieck, Jean-Louis Verdier, Nicolas Bourbaki, Pierre Deligne, Bernard Saint-Donat, Michèle Raynaud, Michel Raynaud, Nicholas Katz, Michel Demazure et Luc Illusie.\par\medskip
Les crédits des autres volumes de la collection, conservés ici dans leur portée propre, nomment : Bas Edixhoven; Yves Laszlo; Michel Raynaud; L. Bayle; N. Borne; O. Brinon; J. Buresi; M. Chardin; F. Ducrot; P. Graftiaux; F. Han; P. Karwasz; L. Koelblen; D. Madore; S. Morel; D. Naie; B. Osserman; J. Riou; V. Sécherre; C. Sabbah; Jacob C. Reinhold; Philippe Gille; Patrick Polo; Daniel Miller; Société mathématique de France.\par\medskip
Travail d'IA consigné : OpenAI Codex (GPT-5.6 Sol, traduction ; GPT-6 Astra, corrections et assemblage), Claude Opus 5.5 (Anthropic), GPT-6.1 Sol (OpenAI API). Le contrôle des scans du 11 novembre 2026 a été effectué par Claude Opus 5.5 et GPT-6.1 Sol. La transcription LaTeX humaine d'Orgogozo reste créditée comme telle ; les couches produites par IA portent les mentions AI-transcribed, AI-translated et AI-produced corrections.\par\medskip
La présente construction est due à GPT-6 Astra (OpenAI, ChatGPT Pro mode), le 12 novembre 2026. La traduction anglaise suit le dépôt sga-en, état 4a4cbac, et son édition du 12 novembre 2026.\par\medskip
Les droits sur l'œuvre originale appartiennent à son auteur, à ses représentants ou à ses ayants droit. Dans la mesure où nous détenons des droits sur la transcription, la traduction, les corrections et la composition produites par l'IA, nous les consacrons au domaine public sous CC0 1.0. Cette édition est publiée à des fins éducatives et non lucratives, et nous ne revendiquons aucun droit que nous ne détenons pas. Si vous détenez des droits sur cette œuvre et vous opposez à sa publication ici, veuillez ouvrir une issue GitHub et nous la retirerons.\par
\medskip\noindent\url{https://github.com/KokunoYumeto/sga-en/issues}
\par\noindent\url{https://creativecommons.org/publicdomain/zero/1.0/}
\par\endgroup\clearpage
